\chapter{Buildings, cities, urbanism}
\label{vol12:ch07}

\epigraph{The city, however small, is in fact divided into two, one the city of the poor, the other of the rich; these are at war with one another.}{Plato, \emph{Republic}, Book IV}

\chapmeta{Half-life: medium-life. Archetypes: scaling, ethics, balance. Exercise target: 25. Project track: design (a $10$-year plan for one specific neighbourhood). Hazard class: medium (multi-decade lock-in; equity stakes). Reader path default: standard, with the building-system and urban-systems sections core and the urban-renewal failure section core. Named cases: Pruitt-Igoe demolition St.\ Louis 1972 (\cite{hist:bristol-pruitt-igoe-1991}); Grenfell Tower fire London 2017 (\cite{acc:grenfell-inquiry-2024}); Robert Moses-era urban renewal NYC (\cite{hist:caro-powerbroker-1974}).}

\noindent
Buildings are the largest engineered artifacts most humans live inside; cities are the largest engineered systems-of-systems most humans live inside; urbanism is the disciplined design of cities under the joint constraints of land, capital, energy, transport, demography, climate, and politics. The chapter takes the building as the unit of analysis at one end, the city at the other, and the streetscape-and-neighbourhood as the recurring middle. Globally, buildings account for approximately $40\%$ of energy use and $33\%$ of CO$_2$ emissions (including embodied carbon in materials; \cite{web:globalabc-status-2023}); cities house approximately $56\%$ of the global population, with the urban-population fraction projected to reach $68\%$ by 2050 (\cite{web:un-wup-2018}). The chapter is the engineering treatment of what humans live inside.

\section{The building as engineered system}\pathtag{core}

A building integrates structure (loads, lateral-force resistance, foundations), envelope (thermal, moisture, vapor, air infiltration, daylight, acoustics), HVAC (heating, ventilation, air conditioning), plumbing (water supply, drainage, fire-suppression), electrical (power, lighting, low-voltage systems), and increasingly digital systems (building-automation, security, communications). Each subsystem has its own engineering discipline; the integration is performed by the architect-and-engineer team across the design process. Volume XI Chapter 3 introduced the building case study; this chapter elaborates the civilization-scale implications.

The building's operational energy use is the dominant component of its lifetime energy footprint for most building types in most climates. For a typical office building, operational energy is $50$-$100\,\text{kWh/m}^2/\text{yr}$ depending on climate, building age, and operational profile; for a passive-house residential building, $\approx 15\,\text{kWh/m}^2/\text{yr}$; for an older inefficient building, $200$-$400\,\text{kWh/m}^2/\text{yr}$. The embodied-energy and embodied-carbon components, historically a smaller fraction, are becoming proportionally more important as operational energy drops. For a high-performance new construction, embodied carbon can equal $20$-$40$ years of operational carbon over a $50$-year service life.

\begin{archetype}[Scaling: building energy efficiency]
A typical $1{,}000$-m$^2$ office at $80\,\text{kWh/m}^2/\text{yr}$ uses $80\,\text{MWh/yr}$. Retrofit to a $30\,\text{kWh/m}^2/\text{yr}$ deep-retrofit standard reduces use to $30\,\text{MWh/yr}$, saving $50\,\text{MWh/yr}$. At \$0.10/kWh, the annual savings are \$5{,}000; the typical retrofit capital cost is \$200-\$500/m$^2$ so \$200{,}000-\$500{,}000 for the building, giving a simple payback of $40$-$100$ years that is longer than typical commercial-investment horizons. The engineering economics is the source of the under-investment in retrofits; the public-policy response is the building-code and incentive structures that internalise the carbon-and-energy externality.
\end{archetype}

\begin{estimation}
\textbf{Question.} Estimate the annual investment required to retrofit the existing building stock of a developed country to net-zero-operational-carbon by 2050.

\textbf{Reasoning.} Developed-country built environment: $\approx 30\,\text{m}^2$ of building floor area per capita, so for a country of $10^8$ population, $\approx 3 \times 10^9\,\text{m}^2$. Retrofit cost: $\approx \$500$/m$^2$ average across the stock (some buildings more, some less). Total: $\approx \$1.5$ trillion. Over $25$ years (2025-2050) at uniform pace: $\$60$ billion/yr. For a $\$1$ trillion GDP, $6\%$. For a $\$5$ trillion GDP economy, $1\%$. Conclusion: the building-stock retrofit is among the largest engineering and capital programmes the chapter discusses; achievable at fractions of GDP typical for major industrial transitions but requiring sustained policy support over multiple decades.
\end{estimation}

\section{Urban systems: streets, blocks, neighbourhoods, districts, cities}\pathtag{core}

A city is an emergent ordering of streets, blocks, neighbourhoods, districts. The street is the public-space organising unit; the block is the parcel-organising unit; the neighbourhood is the social-and-service organising unit; the district is the functional-and-political organising unit; the city is the legal-and-economic organising unit. The engineering relevance is at each scale: streets are designed (the engineering scope of Volume XII Chapter 4); blocks are platted and developed (the engineering scope of civil and site-development engineering); neighbourhoods are planned (the engineering scope of community planning and urban design); districts are infrastructured (the engineering scope of utility and service systems); cities are governed (the engineering scope of municipal engineering).

The city's physical form is determined by these layered decisions over decades and centuries. The pre-industrial European city (Paris within the boulevards; Florence; Vienna within the Ring) is the pre-elevator high-density city. The 19th-century industrial city (Manchester; Chicago; New York's tenements) added the railway and the streetcar. The early-20th-century city (the streetcar suburbs; the City Beautiful movement; the garden-city movement) added the new transit modes. The post-WWII car-centric suburb (Levittown; the Australian suburb; the gated-community development) replaced the public realm with the private parcel. The contemporary mixed-use transit-oriented development is the partial reversal. Each layer leaves its physical and infrastructural inheritance; the contemporary engineer works within the legacy of all of them.

\begin{principle}[Path dependence in urban form]
Urban form is path-dependent: the streets, the blocks, the underground utilities, the major buildings persist for a century or more after their construction. A street platted in 1850 is typically still a street in 2050; a sewer laid in 1900 is typically still a sewer in 2050; a building constructed in 1920 is often still a building in 2050. The engineering implication: today's design decisions are determining the urban form of 2100 and beyond. The discipline required: long-horizon thinking that explicitly acknowledges the durability of the decisions being made.
\end{principle}

\section{Density, sprawl, and their engineering implications}\pathtag{standard}

Urban density (population per unit area) determines the engineering economics of nearly every urban infrastructure. Higher density: per-capita land area lower; per-capita road length lower; per-capita utility-pipe length lower; transit-and-walking mode-share higher; per-capita transport energy lower; per-capita building energy often lower (shared walls; smaller envelope-to-floor-area ratio); per-capita ecological footprint lower. Lower density: more living space per person; less crowding; more private green space; lower property values per square metre; greater car dependence and the associated time-and-emissions costs.

The optimal density depends on the function and the location. The dense European city centre at $\approx 100$-$300$ persons/ha; the typical US suburban density at $\approx 10$-$50$ persons/ha; the rural settlement at $\approx 0.1$-$1$ persons/ha. The contemporary planning practice in most growing cities (Tokyo; Singapore; the European capitals; the leading US cities; the major Australian cities) is for density increases in the central and transit-accessible areas while leaving the suburban and rural patterns more stable. The engineering case for the densification: lower per-capita energy, lower per-capita transport, lower per-capita ecological footprint, higher economic productivity. The political-economic friction with the densification: the existing-resident reaction to changes in their neighbourhood; the housing-affordability dynamics that densification both helps (through supply expansion) and complicates (through value capture and displacement).

\begin{archetype}[Balance: density choices]
Higher density delivers lower per-capita resource and emissions intensity at the cost of higher per-capita land cost and (often) lower per-capita private space. Lower density delivers higher private space and lower nominal land cost per unit area at the cost of higher per-capita resource intensity, higher car dependence, and a substantial public-and-private cost in infrastructure and transport. The optimal density is region- and function-specific; the engineering analysis is the trade-off across the multiple objectives. The civilizational implication: a high-density future is more compatible with planetary boundaries than a low-density future; the dispersed automobile-suburban pattern that the post-WWII US developed is unlikely to be extended to the global $10$-billion population.
\end{archetype}

\section{Climate-resilient cities}\pathtag{core}

Climate change is changing the engineering basis of urban design. The principal urban climate exposures: heat (urban heat island; extreme heat days; heat-related mortality and morbidity); flooding (extreme precipitation; coastal storm surge; sea-level rise); drought (water-supply stress; vegetation loss); wildfire (wildland-urban interface; smoke exposure); air quality (PM2.5 from combustion and wildfire; ground-level ozone); coastal erosion. Each exposure has an engineering response at the city scale.

The urban-heat-island engineering response combines: high-albedo roofing and pavement; tree-canopy expansion; green roofs; water features and evaporative cooling; reduction of heat-rejecting equipment; cooling-centre infrastructure for heatwave response. A typical city can reduce ambient summer temperature by $1$-$3\degC$ through coordinated heat-island-mitigation investment; the health-and-energy benefits are substantial.

The urban-flooding engineering response combines: stormwater-management investment (separation of combined sewers; green-infrastructure for runoff reduction; storage for combined-sewer overflow mitigation; coastal-flooding barriers); land-use planning (avoidance of flood-prone development; managed retreat from the highest-risk areas); building-design adaptation (elevation; flood-proofing; resilient construction). The engineering challenge: the design basis is moving as climate changes, and the systems being designed today must work in the climate of 2050-2100. The contemporary practice uses climate-projection ensembles to size new infrastructure; the challenge is the existing infrastructure that was designed for a climate that no longer obtains.

\begin{principle}[Climate-resilient urbanism as engineering scope]
Climate-resilient urbanism is engineering at the city scale that explicitly designs for the climate of the multi-decade future rather than the climate of the recent past. The engineering elements include: revised design basis (sizing infrastructure for projected rather than historical extremes); managed retreat where appropriate (relocating the highest-risk development inland or up); green-infrastructure investment for multifunctional benefits (heat mitigation, stormwater management, ecological connectivity); building-code adaptation. The institutional element: the planning and zoning structures that internalise the climate-risk information in the development process. The engineering profession's contribution is the technical analysis; the political-economic execution is the binding constraint.
\end{principle}

\section{Affordable-housing engineering}\pathtag{standard}

Affordable housing is the engineering, finance, and policy of delivering housing at a cost that the lower-and-middle-income population can afford. The engineering scope includes: efficient building design and construction (typological optimisation; standardisation; modular and panelised construction; mass timber for mid-rise); infrastructure-efficient siting and parcel design; operating-cost minimisation (energy-and-water efficiency; durability and low-maintenance design); regulatory-and-permitting efficiency (engineering response to the cost of regulatory friction). The engineering frontier in the contemporary affordable-housing sector includes mass-timber mid-rise (Brock Commons at UBC; the Mj{\o}st{\aa}rnet in Norway as the early reference; the contemporary multi-jurisdiction expansion under the revised IBC), modular factory-built housing (the Murano in Singapore; the Y:Cube; the various US modular-housing manufacturers), and prefabricated panelised construction at scale.

The engineering is not the binding constraint in most developed-world affordable-housing crises; the binding constraints are zoning-and-land-use restrictions, financing structures, and political-economic friction with neighbourhood-scale densification. The engineering profession's contribution is the credible cost-and-quality demonstration that the housing typologies the policy debate is about are achievable.

\section{The city as institution: planning, zoning, governance}\pathtag{standard}

A city is also a political-and-institutional entity. Planning (the multi-year strategic-direction setting of the city); zoning (the legal-and-regulatory framework that determines what can be built where); governance (the decision-and-accountability structures by which the city is run). The engineering profession's interface with these institutions is substantial: traffic engineers and street designers are part of the planning departments; building engineers interface with code-and-permitting offices; municipal engineers run the operating departments (water, wastewater, transportation, parks). The institutional engagement is the professional engineer's daily practice in most US and many global jurisdictions.

The governance question is whether the institutions are adequate to the engineering and ethical challenges the chapter has discussed. The Robert Moses-era New York (\cite{hist:caro-powerbroker-1974}) is the case of the engineering-and-planning institution operating without the democratic-accountability constraints that subsequent reforms imposed; the post-1968 New York and the contemporary practice in most US cities operates under substantially more accountability, with both the costs (process friction; cost growth; project delays) and the benefits (better stakeholder engagement; more equitable outcomes). The engineering profession's responsibility includes the engagement with the institutional structures that determine whether the engineering recommendations become reality.

\section{Failure: housing crises, urban-renewal disasters, planning failures}\pathtag{core}

The Pruitt-Igoe public-housing complex in St.\ Louis (\cite{hist:bristol-pruitt-igoe-1991}) is the iconic post-WWII public-housing failure: $33$ high-rise buildings constructed 1954-1956; demolished 1972-1976; the failure attributed variously to architectural design, racial segregation, urban policy, and the wider context of the St.\ Louis industrial decline. The engineering interpretation: the architectural form (high-density high-rise public housing in superblocks) was substantially borrowed from European modernist precedent (Le Corbusier's Cit{\'e} Radieuse and its derivatives) but transplanted into a US context with different policies for maintenance, screening, social services, and integration; the failure mode was not principally architectural but the combination of architecture and inadequate institutional support, with the racial-segregation politics of the era as the binding context.

The Robert Moses-era urban renewal in New York and elsewhere (\cite{hist:caro-powerbroker-1974}) is the case of engineering-and-planning that proceeded without adequate accountability to the affected populations. The Cross-Bronx Expressway (1948-1972), built across substantial existing neighbourhoods, displaced approximately $60{,}000$ residents and severed the South Bronx in ways the area has not recovered from; similar urban-highway construction across the US and elsewhere displaced more than a million low-income and minority residents in the post-WWII decades. The engineering profession's complicity: the traffic engineers, the highway engineers, the construction engineers participated in the work without adequately surfacing the distributive consequences. The institutional response (the National Environmental Policy Act 1969; the section 4(f) of the Department of Transportation Act; the Title VI civil-rights enforcement) was the regulatory structure created in response.

The Grenfell Tower fire in London (2017; \cite{acc:grenfell-inquiry-2024}) is the contemporary case: a $24$-storey residential tower in a working-class West London neighbourhood; flammable aluminium-composite cladding installed during a 2015-2016 retrofit; a kitchen fire that escaped the originating flat and rapidly spread up the cladding; $72$ deaths; the public inquiry's seven-year process documented the systemic failures of the building-products certification, the regulatory regime, the contractor practices, and the local authority's response. The engineering interpretation: the failure was the certification-and-product-regulation failure (the cladding system had been tested in configurations that did not include the as-installed assembly; the regulatory framework allowed the gap), compounded by the budget-driven retrofit-procurement that selected the lower-cost cladding option. The institutional response: the Building Safety Act 2022 in the UK; the Higher-Risk Buildings register; the construction-products reform; continuing remediation of comparable cladding across the UK building stock.

\begin{principle}[Engineering for the affected, not for the abstract user]
The chapter's failure cases each share a structural pattern: engineering decisions made with reference to the abstract user (the design population; the modal occupant; the typical resident) rather than with reference to the actual affected population, whose interests and constraints differ. The engineering response: explicit engagement with the affected population's perspective in the design process; institutional structures that make this engagement difficult to bypass; professional-ethics codes that frame the responsibility as continuous through the design, construction, and post-occupancy phases.
\end{principle}

\section{Project}\pathtag{core}

\begin{project}[$10$-year neighbourhood plan]
Choose a specific neighbourhood with available data. Document the current state: demographics; housing stock (typology, age, condition); infrastructure (transport, water, wastewater, energy, communications); land use; open space; community facilities; economic profile. Identify the principal engineering and policy issues over a $10$-year horizon: climate exposure; demographic trends; infrastructure renewal needs; housing-supply and affordability dynamics; transport-and-mobility patterns. Develop a $10$-year plan that addresses the issues; specify the engineering interventions, the institutional changes required, the estimated capital and operational cost, the distributive consequences. Document the community-engagement strategy and the political-economic feasibility analysis.
\end{project}

\section{Exercises}

\subsection*{Calculation}\pathtag{core}

\begin{exercise}
A $20{,}000$-m$^2$ office building uses $90\,\text{kWh/m}^2/\text{yr}$ of electricity. Compute the annual electricity demand. Compare to the per-employee energy demand at typical occupancy of $15\,\text{m}^2$/employee.
\end{exercise}

\begin{exercise}
A residential building's heating load is $50\,\text{W/m}^2$ at $-10\degC$ outdoor design temperature with $20\degC$ indoor. Compute the heating-degree-day-based annual heating demand for a city with $3{,}000\,\text{HDD}_{18}$/yr.
\end{exercise}

\begin{exercise}
A city of $1$ million people at $50$ persons/ha density occupies $200\,\text{km}^2$; at $100$ persons/ha, $100\,\text{km}^2$. Compute the difference in road length, water pipe length, and sewer length assuming per-area infrastructure intensity.
\end{exercise}

\begin{exercise}
An urban heat-island reduction of $2\degC$ in a city of $1$ million people reduces cooling-energy demand by approximately $10\%$. Compute the kWh and CO$_2$ savings at a baseline cooling load of $5{,}000\,\text{kWh}$/household/yr and $300\,\text{gCO}_2$/kWh grid.
\end{exercise}

\begin{exercise}
A modular-construction housing project produces $4$-bedroom homes at \$200{,}000 each vs \$350{,}000 for site-built; the modular requires $4$-month delivery vs $14$-month for site-built. Compute the financing-cost saving at $7\%$ construction-loan rate.
\end{exercise}

\subsection*{Derivation}\pathtag{standard}

\begin{exercise}
Derive the relationship between urban density and per-capita infrastructure cost for a city's water-distribution network. Identify the elasticity of per-capita cost with respect to density.
\end{exercise}

\begin{exercise}
Derive the urban-heat-island intensity from surface energy balance: incoming radiation, albedo, evapotranspiration, anthropogenic heat. Identify the dominant terms for a typical city.
\end{exercise}

\begin{exercise}
Derive the optimal building-envelope insulation thickness as a function of energy price, insulation cost, and heating-degree-days.
\end{exercise}

\begin{exercise}
Derive the embodied-carbon-vs-operational-carbon trade-off for a building over its lifecycle; identify the conditions under which embodied carbon dominates.
\end{exercise}

\subsection*{Estimation}\pathtag{standard}

\begin{exercise}
Estimate the global building-sector floor area in 2050 under continued urbanisation and per-capita floor area trends. Identify the implied construction-material demand.
\end{exercise}

\begin{exercise}
Estimate the urban-flooding capital exposure for a coastal city of $5$ million people under $1\,\text{m}$ of sea-level rise and 1-in-100-year storm surge.
\end{exercise}

\begin{exercise}
Estimate the global affordable-housing gap in low- and middle-income countries (people without adequate housing) and the capital required to close it.
\end{exercise}

\begin{exercise}
Estimate the heat-related mortality reduction from a coordinated urban-heat-island mitigation programme in a major city during typical and extreme heat-wave conditions.
\end{exercise}

\subsection*{Simulation}\pathtag{standard}

\begin{exercise}
Simulate the energy-and-cost outcomes of a portfolio of building-retrofit options (envelope; heating system; ventilation) for a typical building. Identify the cost-effective combinations and the deep-retrofit pathway.
\end{exercise}

\begin{exercise}
Simulate urban-flooding response under different storm-water management investment levels and different climate-change scenarios. Identify the configurations that maintain target flood-protection levels.
\end{exercise}

\begin{exercise}
Simulate the housing-supply response under different zoning-and-development policy parameters. Identify the policies that produce target affordability outcomes.
\end{exercise}

\subsection*{Design}\pathtag{standard}

\begin{exercise}
Design a deep-retrofit programme for a $1980$s-era multifamily residential building in a temperate climate. Specify the envelope improvements, the systems replacement, the energy outcome, the cost, the resident-impact mitigation.
\end{exercise}

\begin{exercise}
Design a climate-adaptive streetscape for a $1\,\text{km}$ commercial street in a hot-and-flooding-prone city. Specify the tree-canopy, the green-infrastructure stormwater management, the pavement treatments, the cooling features.
\end{exercise}

\begin{exercise}
Design a transit-oriented mixed-use development for a $10$-ha site adjacent to a new rail station. Specify the building program, the density, the parking and mobility, the open space, the affordability mix, the phasing.
\end{exercise}

\subsection*{Judgement}\pathtag{mastery}

\begin{exercise}
Argue: the global urban-population expansion to $7$ billion by 2050 must be accommodated in higher-density forms; the suburban automobile-oriented pattern cannot be scaled to the global population. Argue the opposite. Identify the empirical evidence and the political-economic considerations.
\end{exercise}

\begin{exercise}
Discuss the Pruitt-Igoe public-housing failure. Identify the engineering, architectural, and institutional contributions to the failure; identify the lessons for contemporary affordable-housing engineering.
\end{exercise}

\begin{exercise}
Argue: the engineering profession should engage with zoning-and-land-use policy because the engineering recommendations cannot be implemented in the absence of supportive regulatory frameworks. Argue the opposite. Identify the implications for engineering-professional practice.
\end{exercise}

\begin{exercise}
Discuss the Grenfell Tower fire as a building-safety and certification-regime failure. Identify the engineering decisions, the regulatory regime, the procurement decisions. Identify the implications for the global building-safety regime.
\end{exercise}

\begin{exercise}
Discuss the engineering profession's role in urban-renewal failures (Cross-Bronx Expressway; the Robert Moses-era projects). Identify the responsibilities the profession bore; identify the institutional structures that have changed since.
\end{exercise}

\begin{exercise}
Argue: the climate-adaptation investment required for major coastal cities (e.g., Miami, Jakarta, Lagos, Alexandria) is comparable in magnitude to the buildings-sector decarbonisation investment; the two should be integrated in city-scale planning. Argue the opposite. Identify the empirical evidence and the political-economic considerations.
\end{exercise}
